% Working analysis draft. No natural inference or method-win claim.
% Build with scripts/build_paper_v26.py and the pinned, unmodified ACL style.
\documentclass[11pt]{article}
\usepackage[review]{acl}
\usepackage{times}
\usepackage{latexsym}
\usepackage[T1]{fontenc}
\usepackage[utf8]{inputenc}
\usepackage{microtype}
\usepackage{booktabs}
\usepackage{amsmath}
\hypersetup{pdfauthor={},pdftitle={When Recency Is Not Currentness},pdfsubject={Development analysis scaffold}}

\title{When Recency Is Not Currentness:\\Auditing Temporal Interpretation in Retrieval-Augmented QA}
\author{Anonymous ACL submission}

\begin{document}
\maketitle
\begin{abstract}
\textbf{Working analysis draft; not submission-ready.}
Temporal question answering requires matching evidence to the requested entity, scope, and time.
Selecting the newest timestamp does not establish that match.
We audit our earlier experimental package and identify several confounded measurements.
Among 35 current-holder questions, nine selected reference claims do not explicitly support the requested current role.
A global date baseline retains the requested subject--relation key for only five questions.
The reported aggregate improvement cannot reach $p<0.05$ under any compatible exact paired McNemar test.
These findings concern one package, with assistant-assisted source review.
We also report a diagnostic design failure before natural-model inference.
The proposed representation described query slots rather than the bindings of fixed source assertions.
Unchanged query references would therefore not test the intended source-binding hypothesis.
An assistant source audit then examined all 60 fixed evidence expansions.
Second review of eight flagged cases retained three partial examples and no clear, complete task-level witness.
This describes the audited packets, not model performance or the broader hypothesis.
We establish neither a new retrieval algorithm nor an answer-level advantage.
\end{abstract}

\section{Introduction}

An older document may describe a continuing state.
A newer document may describe a different population or a past event.
Before suppressing evidence, a temporal system must establish what the question and sources mean together.
This includes distinguishing an event's start from a document's publication and a claim's continuing validity.
These distinctions motivate our evaluation; they are not proposed as new temporal concepts.

Recent methods already address substantial parts of this problem.
MRAG separates question content from temporal constraints and ranks evidence using semantic and temporal signals \citep{zhang2025mrag}.
Re$^3$ combines temporal representations with conflict-aware recency filtering \citep{cao2026re3}.
Its recency rule assumes that the newest conflicting version supersedes earlier versions within credible corpora.
Chronos organizes evolving evidence into an event graph \citep{liu2026chronos}.
Adding temporal fields or a revision prompt therefore does not, alone, establish a new method.

We instead ask what an evaluation actually establishes before claiming better temporal reasoning.
Our earlier package compared lifecycle filtering with flat retrieval on current-holder questions.
Its apparent advantage could reflect reference construction, topical evidence loss, scoring, or temporal reasoning.
We disentangle those explanations through a retrospective self-audit.
We preserve every original reference and report all questions.

The audit motivates a second, prospective question.
When additional evidence arrives, does a model correct an asserted temporal interpretation beyond ordinary re-extraction?
We prepared a bounded natural-question diagnostic, then blocked inference after finding that it measured the wrong object.
The completed contribution in this draft is the package audit.
The source-binding hypothesis and any broader scientific claim remain untested.

\section{Retrospective Package Audit}
\label{sec:audit}

\paragraph{Materials and procedure.}
The audited package contains claims, reference rules, retrieval code, aggregate results, and cached model outputs.
Its natural-prose comparison contains 35 current-holder questions.
None supplies a query date.
We inspect the reference construction and all 35 selected reference claims.
We additionally inspect complete same-key histories for nine flagged questions.
Source judgments are assistant-assisted annotations, not independent human gold.
Exact source positions and hashes support later review.
The audit makes no new model calls and checks no live-world officeholders.

\paragraph{Current support versus latest start.}
The reference rule selects the maximum timestamp within a curated subject--relation group.
Those timestamps encode tenure starts, rather than publication dates.
Table~\ref{tab:support} summarizes the selected claims' released text.
Nine claims do not explicitly assert the requested current role.
Six explicitly describe ended tenures.
The remaining cases involve past wording, another club, or an absent club name.
These observations establish internal task--evidence mismatches.
They do not identify the actual current holder.
Asking for the latest supplied start would define a different, narrower task.

\begin{table}[t]
\centering
\small
\begin{tabular}{lr}
\toprule
Selected reference claim's released text & Count \\
\midrule
Explicit requested current role & 26 \\
Explicit ended tenure & 6 \\
Most recent past role only & 1 \\
Different club named as current & 1 \\
Requested club absent & 1 \\
\midrule
All questions & 35 \\
\bottomrule
\end{tabular}
\caption{Assistant-assisted review of all selected reference claims. These are source-support categories, not external factuality labels.}
\label{tab:support}
\end{table}

\paragraph{A baseline without topical evidence.}
The date baseline sorts the entire corpus by timestamp before retaining five claims.
Exactly five claims have the maximum timestamp.
Every question therefore receives those same claims, although their order can differ.
Only five questions retain their requested subject--relation key.
Enumerating all possible orders recovers all 35 cached responses through exact prompt hashes.
The recovered score matches the original 5/35 count.
Thus, this baseline's failure cannot establish failure of subject-aware temporal selection.

\paragraph{Scoring and incomplete pairing.}
The original scorer accepts normalized reference containment or a shared name token with at least four characters.
Its stricter variant also rejects listed stale names.
These rules are not strict exact match or independently adjudicated identity correctness.
Full-name phrase matching and normalized equality are therefore reported separately in the audit artifact.
Their disagreements cannot automatically be called false positives.

The cache stores 4,995 unique prompt hashes, but neither prompts nor condition labels.
Bounded reconstruction recovers 139 compatible condition records from 175 possible cells.
These records use 130 distinct cache keys; shared prompts are not independent generations.
Another 29 cells remain unrecovered and seven remain condition-ambiguous.
Missing reconstruction does not establish missing generation.
Flat and lifecycle conditions also use different prompt wording.
The original comparison consequently does not isolate retrieval under an identical reader prompt.

\paragraph{What the aggregate improvement permits.}
The original report gives 30/35 lifecycle successes and 25/35 flat successes.
These margins permit six paired contingency tables.
Exact two-sided McNemar $p$-values range from $0.0625$ to $0.3018$ across those tables.
Even the most favorable pairing does not meet $0.05$ under this test.
This result does not establish equivalence or prove that temporal filtering cannot help.
The 19 reconstructed lifecycle--flat pairs contain two wins, zero losses, and 17 ties under the original scorer.
Their exact paired $p$-value is $0.5$.
This selected subset cannot estimate the missing pairs.

\paragraph{Simpler controls and timestamp precision.}
An existing variant removes older matched claims without comparing their values.
Its mask equals the original mask on five of seven inspected datasets, including the natural-prose set.
With supplied subject--relation keys, newest-value selection matches all 35 natural-prose references by construction.
That oracle-conditioned result is not natural retrieval accuracy.
It shows why key identification and temporal selection need separate measurements.

The original ISO parser also loses day information in 151 of 153 complete-date Wikidata records.
Its boundary condition fails before the ISO time separator and falls back to the year.
Eight groups collapse distinct dates and values into one timestamp.
Using supplied calendar dates changes two selected references.
Eighteen partial dates remain partial; we do not invent missing calendar fields.
The prose derivative shares these histories and is not an independent confirmation.
No original references are overwritten, and no corrected-date reader outputs are generated.

\section{A Blocked Diagnostic Design}
\label{sec:diagnostic}

The audit does not reveal whether interpretation revision offers useful headroom on natural questions.
We prepared a development study, but stopped it before natural-model inference.
The following describes its design and the reason for that decision.

\paragraph{Questions and admission.}
TEMPO contains 1,730 naturally occurring Stack Exchange questions across 13 domains \citep{abdallah2026tempo}.
Its documents and assisted annotations support a natural-question diagnostic, rather than a source-revision replay.
Our frozen roster contains 60 original questions, selected through a recorded hash rule and domain round-robin ordering.
Six previously inspected questions are excluded before sampling.
Sampling uses identifiers and domains, without predictions or question contents.
All 60 remain in the accounting, including unanswerable, ineligible, or failed items.
This domain-balanced sample does not estimate the benchmark's natural domain mixture.

Question-only admission checks whether the original request has an objectively reviewable temporal target.
Merely mentioning a duration does not establish such a target.
Broad causal explanations and code review cannot silently become narrower date questions.
We record representation limits, missing anchors, and any dependence on inaccessible images.
These assistant annotations do not enter inference inputs.

\paragraph{Interpretations and evidence.}
An interpretation assigns entity, relation, scope, time, comparison, and replacement.
Unknown values differ from inapplicable fields.
Two apparently conflicting claims may coexist because their scopes or periods differ.
An older timestamp alone does not establish replacement.
This schema is a measurement instrument, not a novel temporal representation.

We constructed immutable packs $C_0\subseteq C_1$ before predictions.
The expanded pack does not depend on an initial model error.
Every unit retains source identifiers, positions, and hashes.
Question dates, publication dates, and omitted premises remain unknown unless supplied.
The pool combines supplied positive and negative document identifiers, with their role labels removed before lexical ranking.
This conditioned pool cannot establish ordinary retrieval performance.
All 60 items remain; one failed token preflight.
These are preparation checks, not interpretation outcomes.
Missing evidence within $C_1$ cannot establish global unanswerability.

\paragraph{Unexecuted controls.}
The design compared initial extraction, ordinary re-extraction, candidate-preserving revision, and fresh extraction without an initial draft.
The three expanded-context arms would receive the same evidence.
The fresh arm would test draft anchoring and retain its cheaper native cost.
References and aliases would be frozen before predictions.
The revision prompt was an exploratory control, not an algorithmic contribution.
No natural correction, harm, recall, or cost comparison was run.

\paragraph{The wrong measurement object.}
The implemented fields chiefly describe what the question requests.
Assistant drafts label 40 items resolved, two insufficient, and 18 ineligible.
All 60 have identical initial and expanded query-frame references.
These drafts lack independent reference review and are not human gold.
For example, the requested entity and interval can remain fixed as additional evidence changes a source assertion's interpretation.
Stable query references therefore do not demonstrate that source bindings never need revision.
Conversely, an improved query-slot prediction would not establish repair of a source claim.
The diagnostic cannot distinguish these claims without identifying the fixed assertion being interpreted.

The intended hypothesis concerns that assertion's entity, scope, temporal validity, and relation to other claims.
Testing it requires tracking the same source assertion through evidence expansion.
Its initial and revised bindings need source-grounded adjudication.
A changed answer, a new fact, or a resolved unknown is not automatically a corrected binding.
We did not retrofit these requirements into the blocked schema.
Its preparation artifacts and original scaffold remain preserved as a design audit.

\section{Completed Source Audit and Decision}
\label{sec:sourceaudit}

\paragraph{Source criterion.}
Assistant reviewers examined all 60 frozen question--packet pairs.
A candidate needed an unchanged, question-relevant assertion in $C_0$ and additional evidence clarifying its binding in $C_1$.
New answer facts, unrelated clarifications, and unresolved contradictions did not qualify.
Review used the delivered packets, without benchmark answers, omitted source passages, or model predictions.
Every item remained in the accounting.

\paragraph{Candidate review.}
The first pass flagged two possible witnesses, items 026 and 032.
It also documented nearby cases rather than discarding them.
Further assistant review found no clear, task-relevant witness requiring added $C_1$ evidence.
Among eight flagged cases, that review classified three as partial and five as unsupported.
The second reviewer had seen first-pass judgments.
In item 026, added evidence genuinely supplies case-scope detail.
However, $C_0$ already qualifies the broader legal statement's authority as obiter.
In item 032, an added passage clarified maritime-register scope.
That clarification did not address the requested civil-certificate question.
Item 038 already distinguished policy alignment from arms purchases within $C_0$.
Item 016 lacked a secure link establishing that the added account clarified the same assertion.

Additional nearby cases concerned an annualization assumption, unresolved address scope, a question's plotted curve, and an unrelated pension rule.
They illustrate different exclusion reasons, rather than measured errors from a model.
The audit retained no clear, complete task-level witness in these fixed packets.
This is a descriptive finding about reviewed candidate cases.
It does not estimate the prevalence of revisable assertions in TEMPO or natural documents generally.

\paragraph{Research decision.}
We stopped this packet-based mechanism study before natural inference.
There is no measured correction rate, baseline failure rate, or model advantage.
The previously proposed 12-error and six-residual feasibility thresholds were never tested.
Stable query-frame drafts do not refute the source-binding hypothesis.
The source audit likewise does not show that models cannot misread complete evidence.
It shows that these selected expansions did not establish the intended added-evidence mechanism.
Any future investigation requires newly justified targets and a separately frozen design.
No further experiment or answer-level claim is reported here.

\section{Conclusion}

Our package audit separates current support from the latest supplied start date.
It exposes topical baseline failure, incomplete pairing, permissive scoring, and timestamp precision loss.
These findings weaken the earlier empirical claim without establishing a replacement method.
We also identified a mismatch between query-slot extraction and the intended source-binding task before natural inference.
The completed source audit did not justify proceeding with these packets.
That stop decision concerns research feasibility, not a demonstrated failure of source-binding revision.
External validation remains necessary before a broader analysis claim is justified.
\label{v26-main-end}

\clearpage
\section*{Limitations}

This is a working analysis draft, not a submission-ready method paper.
The completed audit concerns our own experimental package.
It does not measure the prevalence of these problems across published systems or public benchmarks.
Source-support judgments use assistant assistance and require independent human adjudication before stronger claims.
We have not checked the actual current officeholders or reconstructed missing archived publication dates.

Cache reconstruction is bounded, and the original cache omits condition labels and complete runtime provenance.
The recovered subset is selected by reconstructability.
Original name scoring also cannot establish strict semantic correctness.
The aggregate paired test addresses statistical uncertainty under that original scoring rule.
It cannot repair task-definition or measurement problems.

The blocked natural study is development material, with no held-out confirmation.
Its compact schema cannot represent every natural question.
Exact aliases can reject valid paraphrases, and observed interpretation classes need not exhaust all valid alternatives.
Pack sufficiency does not certify corpus coverage.
Static evidence expansion does not test historical source replacement or adaptive retrieval quality.
The query-frame design cannot establish failure or success of the source-binding hypothesis.
The source audit uses assistant judgments and a capped, benchmark-conditioned development pool.
Candidate review does not provide independent human labels or corpus-wide absence guarantees.
The second review covered eight flagged cases and was not blinded to the first-pass judgments.
The unreviewed query-frame drafts are unsuitable as confirmatory references.
No natural method win has been demonstrated.

\section*{Research Assistance and Data Handling}

AI assistants supported literature review, source annotations, code preparation, and manuscript drafting.
Their annotations are explicitly labelled and are not independent human judgments.
This scaffold contains no personal identity or project repository link.
An anonymous submission supplement remains to be assembled and reviewed separately.
Public project files do not themselves constitute an anonymous supplement.
Third-party source text remains subject to its original rights.
Hashes and offsets support auditing without assuming permission to redistribute complete captures.

\bibliography{interpretation_v26}
\end{document}
